\VerdierPhysicalPageAnchor{0037}%
\VerdierSemanticAnchor{ch1-1-2-1-coherence}
and for every triple $(g,h,k)$ of elements of $G$, the relation

\VerdierSemanticAnchor{ch1-1-2-1-eq3}
\begin{equation}
  \bigl(T(g)\star c(k,h)\bigr)c(kh,g)
  =\bigl(c(h,g)\star T(k)\bigr)c(k,hg).
  \tag{1.2.1.3}
\end{equation}
holds.

\VerdierSemanticAnchor{ch1-1-2-2}
\noindent\textbf{Definition 1.2.2.}\quad
Let $G$ be a group. A $G^\circ$-\emph{category}
($G^\circ$ being the group opposite to $G$) is a category $\mathcal C_e$
equipped with a family of functors $T(g)$ ($g\in G$) and isomorphisms of
functors $c(h,g)$ ($g,h\in G$) satisfying properties (1.2.1.1),
(1.2.1.2), and (1.2.1.3).

\VerdierSemanticAnchor{ch1-1-2-2-transition}
A $G^\circ$-category is therefore a category $\mathcal C_e$ equipped with
an action up to isomorphism of a group $G^\circ$. We use the notation
$(\mathcal C_e,G^\circ,T,c)$ for the $G^\circ$-category; the isomorphisms
$c(h,g)$ are called the \emph{transition isomorphisms}.

\VerdierSemanticAnchor{ch1-1-2-2-strict}
When the transition isomorphisms are identity isomorphisms, which in
particular implies $T(hg)=T(g)T(h)$ for every pair $g,h\in G$, the
$G^\circ$-category is said to be \emph{strict}. Such an object is therefore
specified by a representation of $G^\circ$ in the automorphism group of
the category $\mathcal C_e$.

\VerdierSemanticAnchor{ch1-1-2-3}
\noindent\textbf{1.2.3.}\quad
Now let $\mathcal C$ and $\mathcal C'$ be two $G$-graded categories and let
$F\colon\mathcal C\to\mathcal C'$ be a graded functor. Choose normalized
cleavages on $\mathcal C$ and $\mathcal C'$. The choice of these cleavages
determines isomorphisms of functors

\VerdierSemanticAnchor{ch1-1-2-3-eq1}
\begin{equation}
  m(g)\colon F_eT(g)\longrightarrow T'(g)F_e,\qquad
  g\in G,\qquad m(e)=\text{the identity on }F_e
  \tag{1.2.3.1}
\end{equation}
such that for every pair $g,h\in G$ the following diagram commutes:

\VerdierSemanticAnchor{ch1-1-2-3-eq2}
\begin{equation}
\begin{tikzcd}[column sep=huge,row sep=large]
F_eT(hg)
  \arrow[r,"m(hg)"]
  \arrow[d,"{F_e\star c(h,g)}",swap]
&
T^{\prime}(hg)F_e
  \arrow[dd,"{c^{\prime}(h,g)\star F_e}"]
\\
F_eT(g)T(h)
  \arrow[d,"m(g)\star T(h)",swap]
&
{}
\\
T^{\prime}(g)F_eT(h)
  \arrow[r,"T^{\prime}(g)\star m(h)"]
&
T^{\prime}(g)T^{\prime}(h)F_e\mathpunct{.}
\end{tikzcd}
\tag{1.2.3.2}
\end{equation}

\VerdierSemanticAnchor{ch1-1-2-4}
\noindent\textbf{Definition 1.2.4.}\quad
Let $(\mathcal C_e,G^\circ,T,c)$ and
$(\mathcal C'_e,G^\circ,T',c')$ be two $G^\circ$-categories. A
$G^\circ$-\emph{functor} is the datum of a functor

\VerdierSemanticAnchor{ch1-1-2-4-fe}
\[
  F_e\colon\mathcal C_e\longrightarrow\mathcal C'_e
\]
